% Generated draft: requires mathematical and source review.
\documentclass[11pt]{article}
\usepackage[margin=25mm]{geometry}
\usepackage{fontspec}
\setmainfont{Latin Modern Roman}
\usepackage{amsmath,amssymb,mathtools,unicode-math}
\setmathfont{Latin Modern Math}
\usepackage{enumitem,needspace,xurl,hyperref}
\setlength{\parindent}{0pt}
\setlength{\parskip}{5pt}
\setlength{\emergencystretch}{3em}
\newcommand{\N}{\mathbb{N}}
\newcommand{\Z}{\mathbb{Z}}
\newcommand{\Q}{\mathbb{Q}}
\newcommand{\R}{\mathbb{R}}
\newcommand{\C}{\mathbb{C}}
\newcommand{\F}{\mathbb{F}}
\newcommand{\A}{\mathbb{A}}
\newcommand{\PP}{\mathbb{P}}
\newcommand{\E}{\mathbb{E}}
\newcommand{\eps}{\varepsilon}
\newcommand{\dd}{\,\mathrm{d}}
\begin{document}
\section*{UnsolvedMath ID: 1400}
\textbf{Source problem code:} catalogue entry
\subsection{Definitions and mathematical statement}
Let $G=(V,E)$ be a finite undirected simple graph, put $n=|V|$, and let $\Delta(G)=\max_{v\in V}\deg_G(v)$ be its maximum degree. A proper edge coloring assigns colors to the edges so that edges sharing a vertex have different colors. The chromatic index $\chi'(G)$ is the least number of colors in such a coloring. Vizing's theorem, cited in the supplied background, places $\chi'(G)$ in $\{\Delta(G),\Delta(G)+1\}$; these cases are called class 1 and class 2, respectively.

A subgraph $H\subseteq G$ is \emph{overfull} if
\[
 |E(H)|>\Delta(H)\left\lfloor\frac{|V(H)|}{2}\right\rfloor.
\]
The subgraph condition in the source additionally requires $\Delta(H)=\Delta(G)$. Thus the relevant condition can be written
\[
 \exists H\subseteq G:\quad \Delta(H)=\Delta(G)
 \quad\text{and}\quad
 |E(H)|>\Delta(G)\left\lfloor\frac{|V(H)|}{2}\right\rfloor.
\]
Here the number of vertices in the overfull inequality belongs to $H$, not to the ambient graph unless $H=G$.

\textbf{Imported conjecture.} For every such graph $G$ satisfying $\Delta(G)\geq n/3$, is the equivalence
\[
 \chi'(G)=\Delta(G)+1
 \quad\Longleftrightarrow\quad
 \exists H\subseteq G\text{ with }\Delta(H)=\Delta(G)
 \text{ and }|E(H)|>\Delta(G)\left\lfloor\frac{|V(H)|}{2}\right\rfloor
\]
true? The non-strict threshold $\Delta(G)\geq n/3$ is exactly the threshold in the imported short statement and background.

\textbf{Threshold qualification.} The literature triage dated 2026-08-17 instead describes the Overfull Conjecture and the remaining range using $\Delta(G)>n/3$. The two domains differ when $n$ is divisible by $3$ and $\Delta(G)=n/3$. Accordingly the displayed imported formulation is not silently identified with the strict-threshold literature formulation; the intended treatment of equality requires review.

The supplied background notes that an overfull subgraph with the stated maximum degree forces class 2. The dated review describes the full strict-threshold conjecture as unresolved and reports progress for sufficiently large robust expanders in a high-degree regime, including a stated consequence for sufficiently large graphs with $\Delta(G)\geq(1+\varepsilon)n/2$. It does not reproduce the exact theorem hypotheses, robust-expansion definition, or dependence of the size threshold on $\varepsilon$. These partial-progress descriptions are retained as supplied context without upgrading them into a fully quantified theorem. The registry label \emph{open} and triage label \emph{partially solved} are dated classifications, not an independent literature verification.
\subsection{Short English statement}
For a simple graph of sufficiently large maximum degree relative to its order, the Overfull Conjecture asks whether needing one extra edge color is exactly explained by a subgraph with too many edges for that many matchings. The import uses $\Delta\geq n/3$, whereas its literature review uses $\Delta>n/3$; the equality case must be checked.
\subsection{Sources}
\begin{itemize}
\item UnsolvedMath local registry, data/unsolvedmath/problems.json, permanent ID 1400, source code catalogue entry, ``Overfull Conjecture.'' Its short statement and background use the non-strict $n/3$ threshold.
\item G. Chen, J. McDonald and S. Shan, ``Towards the Overfull Conjecture II,'' arXiv:2607.02270 (2026), \url{https://arxiv.org/abs/2607.02270}. The supplied review dated 2026-08-17 cites it for the strict-threshold open conjecture and robust-expander/high-degree progress. The paper was not independently consulted for this draft.
\item Ulam AI and the UnsolvedMath Contributors, \emph{UnsolvedMath: A Curated Collection of Open Mathematics Problems}. Dataset curation and original metadata are supplied under CC BY 4.0; referenced source material retains its own terms. Vizing's theorem is named in the background without a separate bibliographic citation.
\end{itemize}
\end{document}
